% ===== §4 =====
\section{Relational Value Without a Place to Be Priced}
\label{p26:sec:value}

\noindent
This section takes the third dimension, the value generated within a relation, and establishes that such value cannot be named as a magnitude, because quantification presupposes an external place from which worth is assessed, and a value that exists only within a relation grants no such place. The result is that relational value cannot be extracted or priced, and this is what keeps it living in the bond rather than available to be drawn off. The section proceeds through the phenomenon, a review of the accounts that would render relational value assessable, the isolation of the core, and the claim.

\subsection*{The phenomenon}

Within a relation, value is generated. Something comes to matter, to weigh, to be worth caring for, and this worth is not carried in from outside but arises in the bond itself. Language seems able to reach this value. We have the vocabulary of worth ready: what this means to me, how much you matter, what our years together are worth. These phrases present value as if it were a magnitude, a quantity that speech could name and, naming, fix.

They do not fix it. The value generated in a relation is not a property lodged in either person, to be read off and reported. It is generated in the relation and exists only there, in the between, and this is what defeats the naming. To name a value as a magnitude is to assume a place from which the naming is done, a standpoint outside the thing valued from which its worth can be assessed and set down. For a value that exists only within a relation there is no such outside place.

\subsection*{Two ways of trying to assess it, and why they fail}

Two lines of thought promise that relational worth can be brought into a common measure, and each fails in a way that shows the structure of the limit.

The first treats the value in a bond as a form of capital. Bourdieu extended the economic vocabulary of capital to goods that are not money, cultural capital and social capital among them, and analyzed relationships as holdings that can be accumulated, converted, and made to yield \parencite{bourdieu1984distinction}. On this account the worth generated in a relation would be a stock of social capital, in principle measurable by what it can be converted into. The account captures something real about how relationships function within fields of power, and it fails at exactly the point that concerns us. The value that social capital measures is the value a relationship has \emph{for} something outside it, its convertibility into advantage; the value generated \emph{within} the bond, the worth the bond has in itself and for those in it, is precisely what does not convert, what cannot be cashed into any external advantage without ceasing to be the thing it was. Social capital measures the relation from outside, in the currency of what it yields; the value in question is the one that has no yield outside the relation, and so no place in the ledger of capital.

The second is the appeal to a common measure of value as such, the assumption, built into any optimization or cost-benefit frame, that distinct goods can be ranked on a single scale so that trade-offs among them can be computed. Raz argued against exactly this assumption for certain goods, holding that some values are incommensurable, that there is no common measure by which, for instance, the worth of a friendship and a sum of money could be compared, and that to attempt the comparison is not to make a difficult judgment but to commit a kind of error about what the values are \parencite{raz1986morality}. The worth generated within a relation is incommensurable in this sense. To ask what it is worth in any external currency, to seek its magnitude on a common scale, is not to pose a hard question but to misunderstand the value, which is not the kind of thing that has a magnitude on a scale shared with other things. Raz names the error; the present account gives its ground, in the next paragraph.

\subsection*{The core}

The core language cannot reach here is the \emph{relational generation} of value, its arising in the between rather than in either party, which is also its refusal of any external place from which it could be quantified. The point is not that the value is large, or hard to compute, or private. The point is structural, and it grounds the incommensurability Raz observed. Quantification presupposes an outside knower who could assign a magnitude to any worth, a position from which all value is in principle assessable, a place occupied, in the theological figure the series has used elsewhere, by an Other who could price everything. A value generated only within a relation denies that such a position exists for it, since to occupy the position would be to stand outside the relation, which is to stand where the value is not. There is no assessor clear of the bond from which the bond's own generated worth could be measured, because the worth does not exist anywhere the assessor could stand. Incommensurability, at this level, is not a curious fact about certain goods; it is the consequence of a relational ontology of value, on which there is no external vantage from which relational worth could be brought to a common measure, because relational worth is nowhere but in the relation.

\begin{claim}[No external place to price the relation]
Value generated within a relation cannot be assigned a magnitude, because quantification requires a standpoint outside the valued thing from which worth is assessed, and relational value exists only within the relation, offering no such standpoint. The incommensurability of relational goods is not a brute datum but a consequence of the relational generation of value: there is no Other positioned to price what exists only in the between.
\end{claim}

\begin{casebox}{the years that have no price}
Asked what a long shared life is worth, one finds the question not hard but wrong, and the wrongness is instructive. It is not that the figure is large or hard to compute. It is that there is no currency the worth converts into, no external ledger where it could be entered, because the worth exists only in the bond and nowhere a price could be set. To name a sum would not undervalue the years; it would misdescribe them, treating as a magnitude on a common scale what has its being only in the relation.
\end{casebox}

\subsection*{Why the limit is generative here}

This is what keeps relational value from being a possession or a price. Because the worth generated in a bond cannot be set down as a magnitude, it cannot be extracted from the bond and carried off, cannot be banked, traded, or realized elsewhere. Consider the counter-case, a relational value that \emph{could} be fully priced. Priced, it could be sold; realized in an external currency, it could be drawn off, and the relation would be left holding a figure where the bond had been. The limit forbids this. Because relational value has no external place to be assessed, it has no external place to be cashed, and so it stays where it was generated, in the living bond, unavailable to extraction. This will matter in Part Four, where the capture of relational value by external powers is shown to work only by first supplying the very forms in which the value is generated, since the value itself, once generated within the bond, offers capital no place to seize it. For now the point is only that the worth generated within a relation, having no place to be priced, has no place to be taken, and that this ungraspability is what keeps it a living worth rather than a realizable asset.
